% 【基础正文】所属：第9章（由 chapters/revised/09-comparative.tex \input 引入）
\minitopic{对应问题}“认识论”作为学科分类具有特定历史。不同传统当然都讨论学习、判断、错误、言说和智慧，但不必把这些议题组织为“知识的必要充分条件”。比较的第一步不是问“中国有没有西方式认识论”，而是明确比较目的。 至少有三种工作：\textbf{历史解释}力求说明文本在原语境中的意思和功能；\textbf{比较分析}选择可比问题，揭示相似与差异；\textbf{规范重构}借旧资源发展今天可辩护的理论。重构可以富有哲学价值，却不能冒充作者本人的历史主张。

\minitopic{比较的四项检验}第一，\textbf{语义检验}：关键术语的使用范围是否相近？“知”可能指事实掌握、识人、实践智慧或道德觉悟。第二，\textbf{论证检验}：双方是否给出相同前提与推理，而非只有结论听起来类似？第三，\textbf{功能检验}：概念在何种实践中工作，是法庭证明、科学解释、修身还是政治劝诫？第四，\textbf{历史检验}：所声称的影响是否有传承证据，还是仅为结构相似。 通过这些检验后，比较可以有三种结果：真正相近；局部可译但目标不同；表面相似却在关键位置不可直接互译。第三种不是失败，也不自动推出两个传统整体“不可通约”。只要研究者仍能说明差异、给出理由并让双方接受批评，比较就已经在进行；局部不可译首先是限制比较主张强度的信号，而不是停止比较的结论。

\minitopic{四组问题的重新打开}\minitopic{知与行}王阳明知行合一不是简单说“知识必然导致行动”。其工夫论针对把道德认知与实践割裂\parencite{tiwald2014}。它能挑战把知识仅视为静态命题库存的模型，但不能直接反驳JTB，因为二者对象和目标不同。

\minitopic{蔽与视角}荀子的“解蔽”关心片面执著如何妨碍把握整体\parencite{xunzi2014}；庄子通过视角转换质疑固定尺度\parencite{zhuangzi2009}。二者可与认知偏差、立场认识论和语境讨论对话，但一个主张通过秩序与学习矫正，一个更警惕任何视角自居终局。

\minitopic{名、实与公共理由}名家、墨家和因明传统包含丰富的语言、分类与推理研究\parencite{graham1989,dunne2004}。把它们称作“东方形式逻辑”可能方便入门，却容易用现代体系裁剪材料。更稳妥的做法是比较具体推理规范：何者算理由，类推如何受约束，反例如何改变分类。

\minitopic{工夫与认知能力}工夫论把认识能力视为长期练习、身体习惯、情感与关系共同塑造的结果。它可扩展德性认识论，使“可靠能力”不只是一项个体内部机制。但若抽去伦理和共同体目标，只剩“训练提高准确率”，就失去原理论的厚度。

\minitopic{一套问题驱动的方法}面对任何知识声称，可以依次询问：
\begin{enumerate}
  \item 命题是什么，真值条件是否清楚？
  \item 主体以何种态度接受它，确信程度是否校准？
  \item 根据来自知觉、记忆、推理、证言还是制度记录？
  \item 根据与信念之间是否有适当的“基于”关系？
  \item 有无击败者、近邻错误或Gettier式运气？
  \item 专家、工具和共同体如何分担认识劳动与责任？
  \item 权力是否不当改变了可信度和表达资源？
  \item 结论只是真信念，还是形成了理解与行动能力？
\end{enumerate}
这不是知识的八条件定义，而是一张诊断清单。不同问题会突出不同项目。

\begin{argumentbox}
全书的元论证：有限主体无法脱离知觉、记忆、语言、他人和制度；依赖却不等于放弃理性，因为依赖关系可以被记录、比较和纠错；认识自主因而不是无依赖，而是有能力识别根据、分配信任、响应反证并承认未知。
\end{argumentbox}

\minitopic{结语：知道如何继续追问}认识论没有提供一种永不出错的算法。它提供的是区分能力：把真与合理分开，把理由与原因分开，把个人确信与公共证据分开，把知识与理解分开。掌握这些区分，不会消除不确定性，却能使不确定性变得可说明、可讨论、可修正。 好的认识者并非从不犯错，而是错误能够暴露；并非凡事独立验证，而是知道何时信任、何时追问；并非永远保持怀疑，而是在根据不足时暂停，在根据改善时更新。这也是一本问题驱动教材最终希望训练的能力。

\begin{comparison}
比较认识论最有价值的结果，不是宣布某传统“早已有之”，而是让原本自然化的问题重新可见：为何把知识主体设想成孤立个人？为何把解释与实践排除在知识价值之外？反过来，当代分析工具也能要求比较研究给出更精确的文本、论证和反例。
\end{comparison}

\minitopic{翻译是一项理论选择}把一个词译为“knowledge”“reason”或“virtue”，已经选择了比较方向。同一“知”在不同句中可能需要译成知道、识别、理解或智慧。可靠研究应给出原文位置、语境和替代译法，而不是用统一译词制造理论一致。 不可译并不意味着不可理解。研究者可以描述概念在推论和实践中的角色，再说明目标语言缺少单一对应词。保留音译有时能防止过早同一化，但若不提供使用说明，音译也只是把困难隐藏起来。

\minitopic{避免两种中心主义}一种错误以当代英美问题为唯一尺度，只在其他传统中寻找相似答案；另一种错误则因反对西方中心主义而免除非西方主张的论证检验。真正平等的比较既允许传统改变问题，也允许双方被反例和文本证据批评。 比较还要避免“代表整个传统”。儒家、道家、佛教内部各有历史分歧，所谓“西方认识论”也包含理性主义、经验主义、实用主义、女性主义和自然主义等不同方案。分析单位应尽量具体到文本、论证或实践。

\minitopic{案例工作坊：知行合一}若研究者说“王阳明证明知道$p$就必然做$p$”，先检查“知”是否指对任意事实命题的接受，“行”是否指一次外显行为，以及命题是在描述心理规律还是提出道德工夫标准\parencite{tiwald2014}。若这些问题答案是否定的，把它形式化成 $Kp\rightarrow Ap$ 就会误译。 更好的重构是：在相关道德语境中，脱离实践倾向的口头认可不足以算作完成的知。这个主张可与能力之知、德性认识论和动机内在主义对话，但它并不因此与任一理论相同。比较结论应同时写出对话点与剩余差异。
